% Clean ACM-style research paper template
% Draft the paper directly in this file. Replace placeholders as needed.
% Keep ACM formatting if your target venue requires it.

\documentclass[sigconf]{acmart}

% -------------------- Optional layout / metadata settings --------------------
% Add packages only when you need them. acmart already loads common packages
% such as booktabs. Avoid changing margins, fonts, or spacing for ACM submissions.

% Optional short title/authors if the full text is long.
% \renewcommand{\shorttitle}{Short Paper Title}
% \renewcommand{\shortauthors}{Author One et al.}

% Do not add sample copyright, DOI, ISBN, or conference values.
% For a final ACM submission, insert the exact metadata supplied by the venue.
% Example fields (leave commented during ordinary drafting):
% \setcopyright{acmlicensed}
% \copyrightyear{2026}
% \acmYear{2026}
% \acmDOI{DOI supplied by the publisher}
% \acmConference[Conference Acronym]{Full Conference Name}{Conference Dates}{Location}
% \acmISBN{ISBN supplied by the publisher}

\begin{document}

% -------------------- Title and authors --------------------
\title{Your Paper Title Here}

% Add one author block per author.
\author{Author One}
\affiliation{%
  \institution{Department / Institution}
  \city{City}
  \country{Country}
}
\email{author@example.com}

% Example for another author (copy and edit as needed):
% \author{Author Two}
% \affiliation{%
%   \institution{Department / Institution}
%   \city{City}
%   \country{Country}
% }
% \email{author2@example.com}

% -------------------- Abstract and keywords --------------------
\begin{abstract}
% In about 150--250 words, briefly describe:
% (1) the problem and its importance,
% (2) the scope or approach of the paper,
% (3) the main findings or contributions, and
% (4) the implications or conclusion.
Write your abstract here.
\end{abstract}

\keywords{keyword one, keyword two, keyword three}

% Optional ACM Computing Classification System concepts.
% Add these only if required by your target ACM venue:
% \begin{CCSXML}
% <ccs2012>
%   <concept>
%     <concept_id>REPLACE_WITH_VALID_CONCEPT_ID</concept_id>
%     <concept_desc>Replace with the correct concept description</concept_desc>
%     <concept_significance>500</concept_significance>
%   </concept>
% </ccs2012>
% \end{CCSXML}
% \ccsdesc[500]{Replace with the correct concept description}

\maketitle

% ==================================================================
% Suggested structure for a survey/review paper. Rename, reorder, or
% remove sections to fit the actual paper and the target venue.
% ==================================================================

\section{Introduction}
% Context, motivation, problem statement, scope, objectives, and contributions.
% End with a brief outline of the rest of the paper.

\section{Background and Preliminaries}
% Introduce terminology, core concepts, and technical background needed by readers.

\section{Review Methodology}
% Explain where and how literature was searched, search terms, date range,
% inclusion/exclusion criteria, and how papers were selected or categorized.
% For a non-survey paper, replace this with an appropriate methods section.

\section{Taxonomy and Thematic Review}
% Organize the literature by themes, approaches, architectures, or categories.
% Add subsections as needed, for example:
% \subsection{Theme or Category One}
% \subsection{Theme or Category Two}
% \subsection{Theme or Category Three}

\section{Comparative Analysis}
% Compare representative approaches using consistent criteria.
% Discuss advantages, limitations, assumptions, datasets, and evaluation metrics.

\section{Challenges and Open Research Questions}
% Summarize unresolved problems, gaps in current work, and promising directions.

\section{Discussion}
% Synthesize the findings, practical implications, and lessons across themes.
% Avoid repeating the previous sections paper by paper.

\section{Limitations}
% State limitations of the review or proposed work, where applicable.

\section{Conclusion}
% Summarize the main takeaways and their significance. Avoid introducing new claims.

% Optional: include acknowledgments only when applicable and permitted by the venue.
% \begin{acks}
% We thank ...
% \end{acks}

% Optional appendix. Place appendices before the bibliography if required by the venue.
% \appendix
% \section{Additional Material}
% Add supplementary details here.

% Add references in references.bib and cite them with \cite{citation_key}.
\bibliographystyle{ACM-Reference-Format}
\bibliography{references}

\end{document}
